% !TeX root = ../main.tex
\chapter{研究计划与预期成果}

\section{进度安排}
以下时间为相对月份示例，请根据培养计划填写实际起止日期。
每个阶段应设置可检查的交付物，并为实验迭代和论文修改留出时间。

\begin{table}[htbp]
  \centering
  \caption{研究进度安排示例}
  \label{tab:schedule}
  \begin{tabularx}{\linewidth}{@{}p{24mm}XX@{}}
    \toprule
    时间阶段 & 主要工作 & 阶段交付物 \\
    \midrule
    第1--2个月 & 文献调研、问题定义与需求分析 & 调研报告、问题定义与基线方案 \\
    第3--4个月 & 方法设计与初步验证 & 设计说明、初步实验记录 \\
    第5--7个月 & 原型实现与系统实验 & 原型、实验数据与分析结果 \\
    第8--9个月 & 方法完善、论文撰写与修改 & 完整实验结果、论文初稿及修改稿 \\
    \bottomrule
  \end{tabularx}
\end{table}

\section{预期成果}
根据课题实际情况填写拟形成的方法、原型、数据、研究结论和学位论文等成果，
并为每项成果说明验收方式。涉及论文发表、专利等计划时，应按实际要求和可行性表述，避免承诺尚无法确定的结果。
